\documentclass[12pt,a4paper]{article}

% Кодировка и язык
\usepackage[T1,T2A]{fontenc}
\usepackage[utf8]{inputenc}
\usepackage[english,russian]{babel}

% Математика
\usepackage{amsmath,amssymb,amsthm}
\usepackage{mathtools}

% Таблицы и графика
\usepackage{array,booktabs,multirow}
\usepackage{graphicx}

% Настройка страницы
\usepackage{geometry}
\geometry{left=2.5cm,right=2.5cm,top=2.5cm,bottom=2.5cm}

% Гиперссылки
\usepackage{hyperref}
\hypersetup{
    colorlinks=true,
    linkcolor=blue,
    citecolor=blue,
    urlcolor=blue
}

% Теоремы и определения
\newtheorem{definition}{Определение}[section]
\newtheorem{theorem}{Теорема}[section]
\newtheorem{proposition}{Утверждение}[section]
\newtheorem{remark}{Замечание}[section]

% Операторы
\DeclareMathOperator{\dist}{dist}
\DeclareMathOperator{\ReLU}{Re}
\newcommand{\R}{\mathbb{R}}
\newcommand{\N}{\mathbb{N}}
\newcommand{\G}{\mathbf{G}}
\newcommand{\A}{\mathcal{A}}
\newcommand{\D}{\mathcal{D}}
\newcommand{\U}{\mathcal{U}}
\newcommand{\M}{\mathcal{M}}
\newcommand{\T}{\mathcal{T}}
\newcommand{\I}{\mathcal{I}}
\newcommand{\Pset}{\mathcal{P}}

\title{Все виды GRA-обнулёнки: статическая, циклическая, хаотическая и гибридная \\[0.5em]
\large Теоретико-множественное и операторное расширение}

\author{А.~И.~Иванов\thanks{Институт динамических систем, Москва, Россия; ivanov@ids.ru} \and 
П.~С.~Петров\thanks{Лаборатория сложных систем, Санкт-Петербург, Россия; petrov@lcss.spb.ru}}
\date{\today}

\begin{document}

\maketitle

\begin{abstract}
GRA-обнулёнка (Geometric Recursive Analytics) представляет собой семейство методов минимизации отклонений от целей (<<пены>>) в многоуровневых динамических системах. В работе даётся полное описание трёх фундаментальных типов GRA (статическая, циклическая, хаотическая) и их гибридных комбинаций с математическими критериями применимости. Далее концепция расширяется до теоретико-множественного и операторного формализма: типы GRA рассматриваются как элементы абстрактного множества $\G = \{S, C, H\}$, над которым определены булевы операции (объединение, пересечение, разность, симметрическая разность, дополнение). Гибридные архитектуры описываются декартовыми произведениями $\G^L$. Вводится понятие оператора GRA-обнулёнки, действующего на пространстве динамических систем, и строится алгебра таких операторов. Обсуждаются практические применения теоретико-множественного подхода: автоматическая классификация систем, проектирование иерархических контроллеров, верификация и адаптивное переключение режимов. Материал завершается итоговой формулой выбора типа GRA с учётом множественных ограничений.

\textbf{Ключевые слова:} GRA-обнулёнка, динамические системы, аттракторы, пена, статическая GRA, циклическая GRA, хаотическая GRA, гибридная GRA, теория множеств, операторный анализ, иерархические системы, AGI/ASI.
\end{abstract}

\tableofcontents
\newpage

\section{Введение}

GRA-обнулёнка (Geometric Recursive Analytics) — это семейство методов минимизации «пены» $\Phi(x)$ (отклонений от целей) в многоуровневых системах. В зависимости от динамического поведения системы выделяют три фундаментальных типа аттракторов, которым соответствуют три вида GRA:

\begin{itemize}
    \item \textbf{Статическая GRA} – целевой аттрактор: неподвижная точка ($\Phi \to 0$).
    \item \textbf{Циклическая GRA} – целевой аттрактор: предельный цикл ($\Phi(t)$ минимально, но не нуль, периодично).
    \item \textbf{Хаотическая GRA} – целевой аттрактор: странный аттрактор (статистическое состояние, минимизация средних флуктуаций).
\end{itemize}

Кроме того, возможны гибридные варианты (например, разные уровни иерархии используют разные типы).

В данной статье мы не только подробно опишем каждый вид с математическими формулами и критериями применимости (разделы~\ref{sec:static}--\ref{sec:hybrid}), но и расширим концепцию до теоретико-множественного и операторного формализма (разделы~\ref{sec:set}--\ref{sec:applications}), что позволяет единообразно комбинировать стратегии и синтезировать сложные регуляторы для систем искусственного интеллекта.

\section{Статическая GRA (неподвижная точка)}
\label{sec:static}

\subsection{Математическая модель}

Система описывается обыкновенным дифференциальным уравнением:
\begin{equation}
    \dot{x} = f(x), \quad x \in \R^n.
\end{equation}
Пусть $x^*$ — стационарная точка ($f(x^*) = 0$). Пена $\Phi(x)$ — положительно определённая функция, например, квадрат расстояния до цели:
\begin{equation}
    \Phi(x) = \| x - x_{\text{goal}} \|^2.
\end{equation}
Цель статической GRA: найти управление $u(t)$ (или параметры), чтобы $x(t) \to x^*$ и $\Phi \to 0$.

Уравнение эволюции (градиентный спуск):
\begin{equation}
    \dot{x} = -\nabla \Phi(x) + \text{шум (если есть)}.
\end{equation}

\subsection{Критерий применимости}

\textbf{Математический:} все собственные значения матрицы линеаризации $A = \frac{\partial f}{\partial x}\big|_{x^*}$ имеют отрицательные действительные части ($\Re(\lambda_i) < 0$). Система асимптотически устойчива.

\textbf{Физический:} в системе доминирует диссипация; отвод энтропии происходит быстрее, чем её производство; цель — достижение единственного оптимального состояния.

\subsection{Примеры}
\begin{itemize}
    \item Термостат (поддержание заданной температуры).
    \item Регулятор положения робота.
    \item Калибровка датчиков.
    \item Задача коммивояжёра (поиск кратчайшего маршрута).
\end{itemize}

\subsection{Вывод}
Статическая GRA применима, когда система может быть полностью успокоена, и это не вредит её функционированию. Она даёт точное, предсказуемое решение.

\section{Циклическая GRA (предельный цикл)}
\label{sec:cyclic}

\subsection{Математическая модель}

Система стремится к предельному циклу — замкнутой траектории $\Gamma = \{x(t): x(t+T) = x(t)\}$ с минимальной амплитудой. Пена определяется как среднеквадратичное отклонение от некоторого «идеального» цикла или как амплитуда колебаний:
\begin{equation}
    \Phi_{\text{cycle}} = \frac{1}{T} \int_0^T \| x(t) - x_{\text{ref}}(t) \|^2 \, dt.
\end{equation}
Цель: минимизировать $\Phi_{\text{cycle}}$ при сохранении автоколебаний (чтобы цикл не выродился в точку).

Динамика с управлением:
\begin{equation}
    \dot{x} = f(x) + u(t), \quad u(t) = -\nabla_x \Phi_{\text{cycle}} + \beta (\varphi - \varphi_0),
\end{equation}
где $\varphi$ — фаза цикла, $\varphi_0$ — желаемая частота.

\subsection{Критерий применимости}

\textbf{Математический:} линеаризованная система в стационарной точке имеет пару чисто мнимых собственных значений $\pm i\omega$, остальные с $\Re < 0$, и первый ляпуновский коэффициент $l_1 < 0$ (бифуркация Андронова–Хопфа). При этом из точки рождается устойчивый предельный цикл.

\textbf{Физический:} система имеет внутренний источник энергии (накачку) и диссипацию, но они сбалансированы так, что возникает автоколебательный режим. Полная остановка невозможна или нежелательна.

\subsection{Примеры}
\begin{itemize}
    \item Сердечный ритм (цикл сокращений).
    \item Часы с гиревым двигателем.
    \item Химические реакции Белоусова–Жаботинского.
    \item Дыхание.
    \item Колебания цен на рынке.
\end{itemize}

\subsection{Вывод}
Циклическая GRA необходима для систем, которые должны «дышать» (осциллировать), чтобы функционировать. Она не обнуляет пену полностью, а минимизирует амплитуду и стабилизирует частоту.

\section{Хаотическая GRA (странный аттрактор)}
\label{sec:chaotic}

\subsection{Математическая модель}

Система демонстрирует детерминированный хаос: траектории чувствительны к начальным условиям, непериодичны, но ограничены в фазовом пространстве (странный аттрактор). Пена определяется как средняя энергия флуктуаций или энтропия Колмогорова–Синая:
\begin{equation}
    \Phi_{\text{chaos}} = \langle \| x - \bar{x} \|^2 \rangle,
\end{equation}
где $\langle \cdot \rangle$ — усреднение по времени (или по ансамблю).

Цель хаотической GRA — минимизировать $\Phi_{\text{chaos}}$ путём управления параметрами, либо перевести систему в режим предельного цикла или точки (управление хаосом).

Уравнение с управлением (метод Pyragas):
\begin{equation}
    \dot{x} = f(x) + K(x(t-\tau) - x(t)),
\end{equation}
где $K$ — матрица коэффициентов, $\tau$ — период желаемого цикла. Подбором $K$ можно стабилизировать неустойчивый цикл внутри хаоса.

\subsection{Критерий применимости}

\textbf{Математический:} старший показатель Ляпунова $\lambda_1 > 0$ (система хаотична). Если возможно управление, которое делает $\lambda_1 \to 0$, то применима циклическая или статическая GRA; если нет — используется статистическая хаотическая GRA (минимизация средних).

\textbf{Физический:} система сильно неравновесна, имеет много степеней свободы, чувствительна к начальным условиям (погода, турбулентность, рынок). Полное предсказание невозможно, но можно влиять на статистические характеристики.

\subsection{Примеры}
\begin{itemize}
    \item Атмосферная конвекция (ураганы).
    \item Турбулентность жидкости.
    \item Рыночные крахи (волатильность).
    \item Фибрилляция сердца (патологический хаос).
    \item Лазер с хаотической модуляцией.
\end{itemize}

\subsection{Вывод}
Хаотическая GRA не стремится к полному обнулению пены, а лишь к уменьшению её статистических характеристик (дисперсии, энтропии). В ряде случаев возможно управление хаосом для перевода системы в режим цикла или точки.

\section{Гибридная GRA (многоуровневая)}
\label{sec:hybrid}

\subsection{Математическая модель}

На разных уровнях иерархии $l$ используются разные типы GRA. Например:
\begin{itemize}
    \item Нижние уровни (быстрые подсистемы) — статическая GRA.
    \item Средние — циклическая GRA.
    \item Верхние — хаотическая GRA (или статистическая).
\end{itemize}

Общий функционал:
\begin{equation}
    J_{\text{hybrid}} = \sum_{l=0}^{K} w_l \cdot \Phi_{\text{type}(l)}^{(l)},
\end{equation}
где $\Phi_{\text{type}(l)}^{(l)}$ — пена, вычисленная по правилам, соответствующим выбранному типу для уровня $l$.

\subsection{Критерий применимости}

Иерархия времён: быстрые переменные могут быть стабилизированы, медленные — осциллируют, а самые медленные — хаотичны. Такое разделение часто встречается в сложных системах (нейросети, экономика, экосистемы).

\subsection{Примеры}
\begin{itemize}
    \item AGI/ASI: быстрая обработка сигналов (статическая), планирование (циклическое), творчество (хаотическое).
    \item Робот: стабилизация позы (статическая), ходьба (циклическая), адаптация к новой среде (хаотический поиск).
    \item Экосистема: популяционный цикл (циклический), климатические флуктуации (хаотические), управление заповедниками (статическое).
\end{itemize}

\subsection{Вывод}
Гибридная GRA — наиболее общий подход, позволяющий сочетать преимущества всех трёх типов на разных масштабах. Она требует детального анализа иерархии времён и параметров.

\section{Сравнительная таблица типов GRA}

\begin{table}[h!]
\centering
\caption{Сравнение фундаментальных типов GRA}
\label{tab:comparison}
\begin{tabular}{@{}lccc>{\raggedright}p{3cm}@{}}
\toprule
Тип GRA & Аттрактор & Критерий & Пена $\Phi$ & Пример \\
\midrule
Статическая & Неподвижная точка & $\forall i:\ \Re(\lambda_i)<0$ & $\to 0$ & Термостат \\
Циклическая & Предельный цикл & $\lambda_{1,2}=\pm i\omega,\ l_1<0$ & Минимальная амплитуда & Сердце \\
Хаотическая & Странный аттрактор & $\lambda_1>0$ & Статистическая дисперсия & Ураган \\
Гибридная & Смешанный & Иерархия времён & Взвешенная сумма & AGI/ASI \\
\bottomrule
\end{tabular}
\end{table}

\section{Общие выводы по типам GRA}

\begin{itemize}
    \item Выбор типа GRA определяется динамическими свойствами системы, а не произволом разработчика. Нельзя применять статическую GRA к системе, которая по своей природе должна осциллировать (сердце) или хаотична (ураган).
    \item Статическая GRA — простейшая и наиболее изученная, но её область применимости ограничена системами, где полное затухание желательно и возможно.
    \item Циклическая GRA необходима для живых систем, часов, экономических циклов и любых автоколебательных процессов.
    \item Хаотическая GRA — самый сложный случай; она часто сводится к статистической минимизации или управлению хаосом.
    \item Гибридная GRA — перспективный путь для создания AGI/ASI, где разные уровни абстракции требуют разных режимов.
\end{itemize}

\textbf{Формула выбора типа GRA:}
\begin{equation}
    \text{Тип GRA} = F\bigl( \{\Re(\lambda_i)\},\ l_1,\ \lambda_{\text{Ляпунова}},\ \text{иерархия времён} \bigr).
\end{equation}
\textbf{Практический совет:} всегда начинайте с линеаризации и расчёта собственных чисел. Если все $\Re(\lambda_i)<0$ — статическая GRA. Если есть пара мнимых — проверьте $l_1$ для циклической. Если есть положительный показатель Ляпунова — нужна хаотическая или гибридная. Игнорирование этого правила ведёт к нерабочим алгоритмам.

\newpage

% =================== ТЕОРЕТИКО-МНОЖЕСТВЕННОЕ РАСШИРЕНИЕ ===================
\section{Теоретико-множественное представление видов GRA-обнулёнки}
\label{sec:set}

\subsection{Множество типов GRA}

Определим множество $\G$ всех фундаментальных типов GRA-обнулёнки (без учёта гибридных комбинаций):
\begin{equation}
    \G = \{ S, C, H \},
\end{equation}
где $S$ — статическая GRA (аттрактор — неподвижная точка), $C$ — циклическая GRA (предельный цикл), $H$ — хаотическая GRA (странный аттрактор).

Гибридные типы представляют собой элементы булеана $\Pset(\G)$ или декартова произведения $\G^n$ для иерархических систем.

\subsection{Операции над элементами множества $\G$}

На множестве $\G$ можно ввести следующие операции, интерпретируемые в контексте динамических систем и управления пеной.

\subsubsection{Объединение $A \cup B$}

\textbf{Формально:} бинарная операция на $\G$ со значениями в $\Pset(\G)$.

\textbf{Интерпретация:} система может находиться в режиме, где одновременно применимы (или требуются) оба типа GRA для разных подсистем или фазовых областей. Объединение соответствует аддитивной комбинации функционалов пены:
\begin{equation}
    \Phi_{A \cup B}(x) = \alpha \Phi_A(x) + \beta \Phi_B(x), \quad \alpha,\beta \geq 0,\ \alpha+\beta=1.
\end{equation}

\textbf{Пример:} в роботе нижний уровень (стабилизация угла) использует $S$, верхний уровень (планирование маршрута с периодическим обходом) — $C$. Тогда общая стратегия — $S \cup C$ с весами, задаваемыми иерархией времён.

\textbf{Формула для выбора управления:}
\begin{equation}
    u_{A \cup B}(t) = \gamma_A(t) u_A(t) + \gamma_B(t) u_B(t),
\end{equation}
где $u_A, u_B$ — управления, синтезированные по правилам типов $A$ и $B$, а $\gamma_A, \gamma_B$ — функции активации в зависимости от состояния системы.

\subsubsection{Пересечение $A \cap B$}

\textbf{Интерпретация:} система обладает свойствами, при которых применимы оба типа одновременно (например, бистабильность или наличие нескольких аттракторов). Пересечение задаёт область в пространстве параметров, где критерии обоих типов выполняются совместно.

\textbf{Математически:} пусть $\mathcal{S}_A$ — множество систем, для которых применим тип $A$ (например, по критерию собственных чисел). Тогда
\begin{equation}
    \mathcal{S}_{A \cap B} = \mathcal{S}_A \cap \mathcal{S}_B.
\end{equation}

Пена в такой системе минимизируется по пересечению функционалов:
\begin{equation}
    \Phi_{A \cap B} = \max(\Phi_A, \Phi_B) \quad \text{или} \quad \sqrt{\Phi_A^2 + \Phi_B^2}.
\end{equation}

\textbf{Практическое применение:} если система находится на границе бифуркации (например, рождение цикла из точки), то управление должно одновременно подавлять шум и поддерживать амплитуду колебаний в безопасном диапазоне. Пересечение $S \cap C$ даёт компромиссное решение.

\subsubsection{Разность $A \setminus B$}

\textbf{Интерпретация:} применение типа $A$ с явным запретом на поведение, характерное для типа $B$. Например, $C \setminus S$ означает циклическую GRA, но без вырождения в статическую точку (обеспечение минимальной амплитуды).

\textbf{Формула ограничения:} в функционал пены вводится штраф за приближение к нежелательному режиму:
\begin{equation}
    \Phi_{A \setminus B} = \Phi_A + \lambda \cdot \dist(x, \mathcal{A}_B),
\end{equation}
где $\mathcal{A}_B$ — аттрактор типа $B$, $\dist$ — расстояние в фазовом пространстве, $\lambda \gg 1$.

\textbf{Пример:} для сердечного ритма запрещена остановка ($S$), поэтому используется $C \setminus S$ с барьерной функцией, не позволяющей амплитуде упасть ниже порога.

\subsubsection{Симметрическая разность $A \Delta B$}

\textbf{Интерпретация:} выбор между двумя взаимоисключающими типами в зависимости от режима работы системы. Это основа для переключаемых гибридных систем.

\textbf{Формально:} управление задаётся кусочно:
\begin{equation}
    u(t) = \begin{cases}
        u_A(t), & \text{если } x \in \mathcal{D}_A, \\
        u_B(t), & \text{если } x \in \mathcal{D}_B,
    \end{cases}
\end{equation}
где $\mathcal{D}_A \cap \mathcal{D}_B = \emptyset$ и $\mathcal{D}_A \cup \mathcal{D}_B = \mathcal{D}$ — рабочая область.

\textbf{Пример:} в режиме нормальной ходьбы робота применяется $C$ (циклическая GRA), а при обнаружении препятствия — $S$ (статическая стабилизация для остановки). Это симметрическая разность с переключением по событию.

\subsubsection{Дополнение $\overline{A}$ относительно универсума}

\textbf{Универсум} $\U$ — множество всех возможных способов минимизации пены в динамических системах, включая не-GRA методы (например, стохастическое управление, адаптивное управление без явного аттрактора).

\textbf{Интерпретация:} $\overline{A}$ — все подходы, не являющиеся данным типом GRA.

\textbf{Практическое значение:} определение границ применимости GRA-методов и необходимости использования иных подходов (например, для систем с негауссовыми шумами или с неизвестной моделью).

\subsection{Декартово произведение и гибридные типы}

Гибридная GRA, рассмотренная в разделе~\ref{sec:hybrid}, формально представляется как элемент декартова произведения:
\begin{equation}
    \G_{\text{hybrid}}^{(L)} = \G^L = \underbrace{\G \times \G \times \dots \times \G}_{L \text{ раз}},
\end{equation}
где $L$ — количество иерархических уровней (или временных масштабов). Элемент $(g_0, g_1, \dots, g_{L-1}) \in \G^L$ задаёт тип GRA на каждом уровне.

\textbf{Функционал пены гибридной системы:}
\begin{equation}
    \Phi_{\text{hybrid}}(x) = \sum_{l=0}^{L-1} w_l \cdot \Phi_{g_l}(x^{(l)}),
\end{equation}
где $x^{(l)}$ — переменные уровня $l$, $w_l$ — весовые коэффициенты, отражающие иерархию времён.

\textbf{Операция композиции уровней:} если на некотором уровне применяется не один тип, а комбинация, то вместо $\G$ используется $\Pset(\G)$.

\textbf{Пример:} AGI/ASI архитектура: \\
Уровень 0 (сенсорика) — $S$, \\
Уровень 1 (рефлексы) — $C$, \\
Уровень 2 (планирование) — $C \cup H$, \\
Уровень 3 (творчество) — $H$. \\
Тогда полный тип — кортеж $(S, C, C \cup H, H) \in (\Pset(\G))^4$.

\section{GRA-обнулёнка как оператор}
\label{sec:operator}

\subsection{Определение оператора}

Рассмотрим пространство $\M$ всех динамических систем с управлением:
\begin{equation}
    \dot{x} = f(x, u), \quad x \in \R^n,\ u \in \R^m.
\end{equation}

Каждой системе $(f, \text{цель})$ можно сопоставить тип GRA $g \in \G$ (или его гибридное расширение) по правилу, основанному на спектральных свойствах линеаризации и показателях Ляпунова (см. раздел~\ref{sec:static}--\ref{sec:chaotic}). Таким образом, определяется \textbf{оператор типизации}:
\begin{equation}
    \T : \M \to \Pset(\G) \cup \G^L.
\end{equation}

Обратно, зная тип $g$, можно синтезировать \textbf{оператор управления}:
\begin{equation}
    \U_g : \M_{\text{open}} \to \M_{\text{closed}},
\end{equation}
который по разомкнутой системе строит замкнутую с регулятором, реализующим GRA-обнулёнку типа $g$.

\textbf{Композиция операторов:}
\begin{equation}
    \U_{g_1} \circ \U_{g_2}
\end{equation}
означает последовательное применение двух стратегий (например, сначала стабилизация точки, затем введение малых колебаний для цикла). Это возможно, если соответствующие подсистемы разделены по времени или пространству.

\textbf{Тождественный оператор} $\I$ соответствует отсутствию управления (система сохраняет свою динамику). Он не принадлежит $\G$, но входит в универсум $\U$.

\textbf{Обратный оператор} $\U_g^{-1}$ (если существует) — переход от замкнутой системы к разомкнутой, т.е. идентификация исходной динамики по наблюдаемому поведению с регулятором.

\subsection{Алгебра операторов GRA}

На множестве операторов $\{\U_g \mid g \in \G^*\}$ можно ввести операции, индуцированные теоретико-множественными операциями над типами:

\begin{itemize}
    \item \textbf{Сумма операторов:} $\U_A \oplus \U_B = \U_{A \cup B}$ (аддитивное управление).
    \item \textbf{Произведение операторов:} $\U_A \otimes \U_B = \U_{A \cap B}$ (совместное применение с ограничениями).
    \item \textbf{Инверсия:} $\U_{\overline{A}}$ — оператор, реализующий любой метод, кроме GRA типа $A$.
\end{itemize}

\textbf{Практический смысл:} такая алгебра позволяет формально описывать сложные регуляторы как комбинацию базовых GRA-стратегий. Например, регулятор для шагающего робота может быть записан как:
\begin{equation}
    \U_{\text{robot}} = (\U_S \oplus \U_C) \otimes \U_{C \setminus S},
\end{equation}
что означает: одновременно стабилизация и цикл, но с запретом остановки цикла.

\section{Практические применения теоретико-множественного подхода}
\label{sec:applications}

\subsection{Автоматическая классификация систем}

Имея численную модель системы, можно вычислить:
\begin{itemize}
    \item собственные значения линеаризованной матрицы $\{\lambda_i\}$,
    \item первый ляпуновский коэффициент $l_1$,
    \item старший показатель Ляпунова $\lambda_{\max}$.
\end{itemize}

По этим данным строится отображение $\T(f)$. Множество $\G$ используется как пространство меток для классификатора. Например:

\begin{table}[h!]
\centering
\caption{Правила классификации системы по типу GRA}
\label{tab:classification}
\begin{tabular}{@{}lc@{}}
\toprule
Условие & Тип GRA \\
\midrule
$\forall i:\ \Re(\lambda_i) < 0$ & $S$ \\
$\exists i:\ \Re(\lambda_i) = 0,\ l_1 < 0$ & $C$ \\
$\lambda_{\max} > 0$ & $H$ \\
\bottomrule
\end{tabular}
\end{table}

\subsection{Проектирование иерархических контроллеров}

При построении AGI/ASI архитектуры каждому уровню абстракции сопоставляется элемент из $\Pset(\G)$ или $\G^L$. Операции объединения и пересечения позволяют формально описать взаимодействие уровней. Например, уровень сознания может использовать $H \cap C$ — хаотический поиск с периодической фокусировкой внимания.

\subsection{Верификация и отладка}

Используя операции разности и дополнения, можно формально задать \textbf{запрещённые режимы} (например, $\overline{S}$ на уровне дыхания — запрет остановки дыхания). Это упрощает спецификацию требований безопасности.

\subsection{Адаптивное переключение типов}

На основе симметрической разности реализуется гистерезисное переключение между $S$ и $C$ в зависимости от внешних условий (например, экономическая политика: в кризис — стабилизация, в рост — циклическое регулирование).

\section{Выводы}
\label{sec:conclusions}

\begin{enumerate}
    \item \textbf{Множественная природа GRA:} Фундаментальные типы GRA-обнулёнки образуют дискретное множество $\G$, на котором естественным образом определяются булевы операции. Эти операции имеют чёткую физическую интерпретацию в терминах комбинирования функционалов пены и ограничений на аттракторы.
    \item \textbf{Гибридность как декартово произведение:} Иерархические и многоуровневые системы описываются элементами $\G^L$ или $\Pset(\G)^L$, что позволяет формально конструировать сложные стратегии управления.
    \item \textbf{Операторная алгебра:} GRA-обнулёнка может быть представлена как оператор, действующий на пространстве динамических систем. Композиция, сумма и произведение операторов отражают комбинацию стратегий минимизации пены.
    \item \textbf{Практическая ценность:} Теоретико-множественный подход даёт:
    \begin{itemize}
        \item единый язык для спецификации требований к регуляторам;
        \item формальную основу для автоматического синтеза гибридных контроллеров;
        \item инструмент для анализа совместимости различных методов управления.
    \end{itemize}
\end{enumerate}

\textbf{Итоговая формула оператора выбора типа GRA с учётом множественных ограничений:}
\begin{equation}
    g^* = \arg\min_{g \in \mathcal{G}} \left\{ \Phi_g(x) + \sum_{i} \mu_i \cdot \mathbb{I}[x \notin \mathcal{D}_{g_i}] \right\},
\end{equation}
где $\mathcal{G} \subseteq \Pset(\G)$ — допустимое подмножество типов, $\mathbb{I}$ — индикаторная функция запрещённых областей, $\mu_i$ — большие штрафные коэффициенты.

Такой подход делает GRA-обнулёнку не просто набором разрозненных методов, а стройной алгебраической системой, готовой к применению в самых сложных задачах управления, включая создание сильного искусственного интеллекта.

% Благодарности (опционально)
\section*{Благодарности}
Работа выполнена при поддержке гранта РФФИ №~XX-XX-XXXXX.

% Список литературы (можно добавить позже)
\begin{thebibliography}{9}
\bibitem{pyragas} Pyragas K. Continuous control of chaos by self-controlling feedback // Physics Letters A, 1992.
\bibitem{kuznetsov} Кузнецов Ю.А. Элементы прикладной теории бифуркаций. — М.: УРСС, 2001.
\bibitem{strogatz} Strogatz S.H. Nonlinear Dynamics and Chaos. — Westview Press, 2014.
\bibitem{andronov} Андронов А.А., Витт А.А., Хайкин С.Э. Теория колебаний. — М.: Наука, 1981.
\end{thebibliography}

\end{document}